\par\Needspace{4\baselineskip}
\begin{center}
\hypertarget{native-061-article044}{}\pdfbookmark[2]{44.}{article-044}44.
\end{center}
We may deal more briefly with \emph{hyperbolic} motion, since it is treated by a method entirely analogous to the one we have described for elliptic motion. We write the equation between the time $t$ and the auxiliary quantity $u$ in the following form.
\[
\begin{aligned}
&(e-1)\left\{\tfrac{1}{20}\left(u-\frac{1}{u}\right)+\tfrac{9}{10}\log u\right\}\\
&\quad+\left(\tfrac{1}{10}+\tfrac{9}{10}e\right)\left\{\tfrac{1}{2}\left(u-\frac{1}{u}\right)-\log u\right\}=kt\left(\frac{e-1}{q}\right)^{\frac{3}{2}},
\end{aligned}
\]
\noindent The logarithms are hyperbolic. Here $\tfrac{1}{20}\left(u-\dfrac{1}{u}\right)+\tfrac{9}{10}\log u$ is a quantity of the first order and $\tfrac{1}{2}\left(u-\dfrac{1}{u}\right)-\log u$ a quantity of the third order, if $\log u$ is regarded as a small quantity of the first order. Thus, on setting
\[
\begin{aligned}
\frac{6\left\{\tfrac{1}{2}\left(u-\frac{1}{u}\right)-\log u\right\}}{\tfrac{1}{20}\left(u-\frac{1}{u}\right)+\tfrac{9}{10}\log u}&=4A,\\[.7em]
\frac{\tfrac{1}{20}\left(u-\frac{1}{u}\right)+\tfrac{9}{10}\log u}{2\BareRoot{A}}&=B,
\end{aligned}
\]
\noindent\SourceLeaf{062}$A$ is a quantity of the second order, and $B$ differs from unity by a quantity of the fourth order. Our equation then takes the following form
\[
\text{2)}\qquad B\left\{2(e-1)A^{\frac{1}{2}}+\tfrac{2}{15}(1+9e)A^{\frac{3}{2}}\right\}=kt\left(\frac{e-1}{q}\right)^{\frac{3}{2}},
\]
\noindent which is entirely analogous to equation 1) of Art. 37. Further, on setting $\left(\dfrac{u-1}{u+1}\right)^2=T$, we find that $T$ is of the second order and, by the method of infinite series, that
\[
\begin{aligned}
\frac{A}{T}&=1+\tfrac{4}{5}A+\tfrac{8}{175}A^2-\tfrac{8}{525}A^3\\
&\quad+\tfrac{1896}{336875}A^4-\tfrac{28744}{13138125}A^5+\text{etc.}
\end{aligned}
\]
\noindent Accordingly, if we set $\dfrac{A}{T}=1+\tfrac{4}{5}A+C$, then $C$ is a quantity of the fourth order, and $A=\dfrac{(1+C)T}{1-\tfrac{4}{5}T}\mathord{\raisebox{.35ex}{.}}$ Finally, the radius vector follows readily from equation VII of Art. 21.
\[
r=\frac{q}{(1-T)\cos\tfrac{1}{2}v^2}=\frac{(1+\tfrac{4}{5}A+C)q}{(1-\tfrac{1}{5}A+C)\cos\tfrac{1}{2}v^2}\mathord{\raisebox{.35ex}{.}}
\]
